\centerline{\bf \Huge Турбо тренировка I}
\centerline{\bf \Huge }

\begin{flushright}
        \scriptsize{На выполнение ровно 2 часа.}
\end{flushright}

\problem{1}
У Олежи есть 4 карточки с числами 0, 1, 2 , 3, а так же 4 карточки с знаками "+" (прибавить), "×" (умножить) и ":" (разделить). Сколькими способами Олежа может составить пример из имеющихся у него карточек? (В примере должны быть использованы все карточки, знаки действий должны стоять между числами и не должно быть деления на 0.)

\problem{2}
Доказать, что не существует таких натуральных $m, n$ что $m^2+n$ и $n^2+m$ являются точными квадратами.

\problem{3}
Может ли число $n!$ делиться на $2^n?$

\problem{4}
Высоты $B B_1$ и $C C_1$ остроугольного треугольника $A B C$ пересекаются в точке $H$. Точка $O-$ центр описанной окружности треугольника $A B C$. Известно, что $\angle B A C=45^{\circ}$. Докажите, что четырёхугольник $O B_1 H C_1$ — параллелограмм.